% Generated draft: requires mathematical and source review.
\documentclass[11pt]{article}
\usepackage[margin=25mm]{geometry}
\usepackage{fontspec}
\setmainfont{Latin Modern Roman}
\usepackage{amsmath,amssymb,mathtools,unicode-math}
\setmathfont{Latin Modern Math}
\usepackage{enumitem,needspace,xurl,hyperref}
\setlength{\parindent}{0pt}
\setlength{\parskip}{5pt}
\setlength{\emergencystretch}{3em}
\newcommand{\N}{\mathbb{N}}
\newcommand{\Z}{\mathbb{Z}}
\newcommand{\Q}{\mathbb{Q}}
\newcommand{\R}{\mathbb{R}}
\newcommand{\C}{\mathbb{C}}
\newcommand{\F}{\mathbb{F}}
\newcommand{\A}{\mathbb{A}}
\newcommand{\PP}{\mathbb{P}}
\newcommand{\E}{\mathbb{E}}
\newcommand{\eps}{\varepsilon}
\newcommand{\dd}{\,\mathrm{d}}
\begin{document}
\section*{UnsolvedMath ID: 800483}
\textbf{Source problem code:} catalogue entry
\subsection{Definitions and mathematical statement}
Write the standard symplectic space as
\[
(\R^{2n},\omega_0),\qquad \omega_0=\sum_{j=1}^{n}\dd q_j\wedge\dd p_j.
\]
The source does not specify the positive integer \(n\). Let \(M\subset\R^{2n}\) be a smooth hypersurface, possibly considered only as a germ. Its characteristic line at \(x\in M\) is
\[
\mathcal L_x=\ker(\omega_0|_{T_xM})
=\{v\in T_xM:\omega_0(v,w)=0\text{ for every }w\in T_xM\}.
\]
A characteristic curve is a connected integral curve of this line distribution. Interpreting ``planar'' in the affine sense, a characteristic curve \(\gamma:I\to M\) is planar if there exist \(a_\gamma\in\R^{2n}\) and a two-dimensional linear subspace \(V_\gamma\subset\R^{2n}\) such that
\[
\gamma(I)\subset a_\gamma+V_\gamma.
\]
The plane may depend on the curve; no common plane or symplectic condition on these planes is stipulated.

\textbf{Problem (Serge Tabachnikov, Problem 5).} Determine how restrictive it is for every characteristic curve on \(M\) to be planar. In particular, describe the constraints this condition places on a smooth hypersurface germ. The source asks an open-ended rigidity or classification question and supplies no specific proposed classification or equivalence relation. For a germ, curves must be understood locally; whether a corresponding global question requires each entire characteristic leaf to lie in one plane is not specified. When \(n=1\), planarity is automatic because the ambient space is already a plane.

\textbf{Supplied background qualification.} The source states, citing only \([\mathrm{KS}]\), that if \(M\) is closed and convex, then \(M\) is an affine symplectic image of a round ball. An affine symplectic map has the form \(x\mapsto Ax+b\), where \(A^*\omega_0=\omega_0\). As written, the background statement has a dimensional mismatch: \(M\) is a hypersurface, whereas a ball is full-dimensional. It may intend the boundary of such an image, or the convex body bounded by \(M\), but neither correction is established by the supplied record. The meaning of ``closed'' and any additional hypotheses of \([\mathrm{KS}]\) also remain unresolved.
\subsection{Short English statement}
Serge Tabachnikov asks what geometric restrictions follow when every characteristic curve of a smooth hypersurface in standard symplectic space lies in an affine plane, allowing the plane to vary with the curve. The question includes hypersurface germs. The ambient dimension, precise classification target, and local versus global scope need clarification. A supplied closed convex rigidity claim refers to an unidentified \([\mathrm{KS}]\) and ambiguously describes a hypersurface as an image of a ball.
\subsection{Sources}
\begin{itemize}
\item Local UnsolvedMath registry, permanent ID \(800483\), exact source problem code \texttt{catalogue entry}; title: \emph{Billiards and Quantitative Symplectic Geometry --- Serge Tabachnikov --- Problem 5}. The registry labels this problem \emph{open}; that label is supplied curation, not an independently verified mathematical assessment.
\item Albach, Barraud et al., \emph{Open Problems in Billiards and Quantitative Symplectic Geometry} (2026), source item ``Serge Tabachnikov --- Problem 5,'' \url{https://arxiv.org/abs/2602.12896}. The record attributes the contribution to Serge Tabachnikov and records extraction by \texttt{arxiv\_latex}, with access date 2026-10-16. These provenance details are supplied by the record; no independent source visit or verification is claimed.
\item The supplied statement cites \([\mathrm{KS}]\) for its closed convex rigidity assertion. No authors, title, publication details, or link identifying \([\mathrm{KS}]\) are supplied, so the cited result and its hypotheses cannot be checked from this record.
\end{itemize}
\end{document}
